\documentclass[12pt]{article}
\usepackage[margin=0.72in]{geometry}
\usepackage{lmodern}
\usepackage{microtype}
\usepackage{multicol}
\usepackage{fancyhdr}
\usepackage{titlesec}
\usepackage{enumitem}
\usepackage{xcolor}
\usepackage[hidelinks]{hyperref}

\setlength{\columnsep}{0.30in}
\setlength{\parindent}{1.05em}
\setlength{\parskip}{0.24em}
\setlength{\headheight}{15pt}
\titleformat{\section}{\large\bfseries\scshape}{}{0pt}{}
\titleformat{\subsection}{\normalsize\bfseries}{}{0pt}{}
\pagestyle{fancy}
\fancyhf{}
\fancyhead[L]{Aquinas on Reason and Will}
\fancyhead[R]{Reading Handout}
\fancyfoot[C]{\thepage}
\renewcommand{\headrulewidth}{0.4pt}

\newcommand{\focusbox}[1]{%
\vspace{0.4em}\noindent\colorbox{gray!12}{\begin{minipage}{0.96\linewidth}#1\end{minipage}}\vspace{0.5em}}

\begin{document}
\begin{center}
{\LARGE\bfseries Are We Free, or Are We Ruled by Fate?}\par\vspace{0.25em}
{\large\scshape Reading 4: Aquinas on Reason, Will, and Law}\par\vspace{0.45em}
\rule{0.82\textwidth}{0.4pt}\par\vspace{0.7em}
\end{center}

\section*{Vocabulary Preview}
\begin{multicols}{2}
\begin{itemize}[leftmargin=*, itemsep=0.18em]
\item \textbf{Human act}: An action that belongs to a person as a rational creature because it proceeds from reason and will.
\item \textbf{Act of a human}: Something that happens in or through a human being, but not as a deliberate moral action.
\item \textbf{Deliberate will}: The power of choosing after reason has considered an end and possible means.
\item \textbf{End}: The goal, purpose, or good for which an action is done.
\item \textbf{Final end}: The ultimate good toward which a person's desires and actions are ordered.
\item \textbf{Voluntary}: Proceeding from an inner principle with some knowledge of the end.
\item \textbf{Violence}: External force against the will.
\item \textbf{Fear}: A pressure that can make an action difficult or mixed, but does not always remove responsibility.
\item \textbf{Concupiscence}: Strong appetite or desire for some perceived good, especially bodily or lower goods.
\item \textbf{Habit}: A stable disposition that makes a person ready to act well or badly.
\item \textbf{Law}: An ordinance of reason for the common good, made by proper authority and made known.
\item \textbf{Natural law}: The rational creature's participation in eternal law; the light of reason by which we know basic goods.
\item \textbf{Practical reason}: Reason used to judge what should be done.
\item \textbf{Common good}: The good of a whole community, not merely private advantage.
\end{itemize}
\end{multicols}

\section*{Introduction}
Thomas Aquinas lived from A.D. 1225 to 1274. He was a medieval Christian theologian and philosopher who drew deeply from Aristotle while also working within Augustine's Christian tradition. His \textit{Summa Theologiae} is written in a question-and-answer form: Aquinas raises objections, gives an answer, and replies to objections.

In this unit, Aristotle gave us the language of voluntary action, Cicero challenged us to think about fate and causes, and Augustine argued that evil comes from the will's free turning toward lower goods. Aquinas now gathers several threads together. He asks what makes an act truly human, how reason and will work together, why habit matters, and how law can guide freedom rather than merely restrain it.

\focusbox{\textbf{Source note:} Adapted and modernized from Thomas Aquinas, \textit{Summa Theologiae}, First Part of the Second Part, selected articles from Questions 1, 6, 49, 90, 91, and 94. This classroom version preserves Aquinas's main sequence for this unit: human acts proceed from reason and will, voluntary action involves knowledge of an end, habit disposes us well or badly, and law is an ordinance of reason ordered to the common good.}

\begin{multicols}{2}
\section*{Reading}

\subsection*{1. What Makes an Action Truly Human?}
Aquinas begins with a distinction. Some actions merely happen in a human being. A person may blink, digest food, stumble, or move absent-mindedly. These are actions \textit{of a human}, but they are not yet human actions in the full moral sense.

Properly human actions are different. They are actions of which a person is master. A person is master of action through reason and will. Therefore, actions are properly called human when they proceed from deliberate will.

This does not mean that every human action is calm or perfectly wise. A person may choose badly, quickly, or under pressure. But if the action proceeds from within the person, and if the person acts with some knowledge of what he is doing and why, then the action belongs to him morally.

Aquinas is sharpening Aristotle's point. To ask whether a person is responsible, we must ask whether the act came from the person as a rational and willing agent, not merely whether something happened through his body.

\subsection*{2. Human Beings Act for an End}
Aquinas says that every properly human action is done for an end. The end is the goal, purpose, or good for which the action is chosen.

The end may be noble or shameful. A student may study in order to learn wisdom, or only in order to be praised. A ruler may seek peace for the common good, or power for himself. A man may give money out of justice, fear, pride, or love. The outward action matters, but the end helps define what kind of action it really is.

For Aquinas, the will is drawn toward what appears good. No one chooses evil simply because it is evil. Even a sinner chooses some apparent good: pleasure, revenge, security, honor, wealth, or control. The moral problem is that the person may choose a lower or partial good against a higher good.

Because human beings act for ends, freedom cannot mean random action without purpose. A free human act is not an explosion without direction. It is an act ordered toward some good that reason presents and the will seeks.

\subsection*{3. The Final End}
Aquinas then asks whether all human actions point toward some final end. He argues that they do. A person may choose many smaller goods, but those smaller goods are chosen because they seem to lead toward happiness, fulfillment, or some complete good.

People disagree about what happiness is. One person treats pleasure as the highest good. Another treats honor, wealth, success, power, knowledge, friendship, or God as highest. But everyone acts under some idea of the good life.

This matters for responsibility. Our choices are not isolated dots. They reveal what we think life is for. A single action may show a deeper order of love and judgment. If a man repeatedly chooses power over justice, he is not merely making separate mistakes; he is being formed toward a false final end.

Aquinas therefore deepens the unit question. We are not only asking, "Did he choose this act?" We are also asking, "What end was he serving, and what kind of person is he becoming by serving it?"

\subsection*{4. What Makes an Act Voluntary?}
Aquinas defines the voluntary as action that comes from an inner principle with knowledge of the end. A stone falls from an inner natural principle, but it does not know where it is going. An animal may pursue food by sense and instinct. A human being can know an end, understand means, deliberate, and choose.

Therefore voluntary action belongs most fully to rational creatures. Human beings can move themselves toward an end because they can understand the end as an end.

Aquinas also says something important about causes. An outside object may attract the appetite. Another person may persuade. God may move all things as First Mover. None of this automatically destroys voluntariness. A voluntary act can have outside influences and still proceed from within the person.

This point prepares students for Hobbes, Hume, and Laplace. Aquinas does not think responsibility requires that an act have no causes at all. The question is whether the act proceeds through the person's own reason and will with knowledge of the end.

\subsection*{5. Violence, Fear, Desire, and Ignorance}
Aquinas then examines what can make an act less voluntary or involuntary.

Violence can be done to the body. Someone can be dragged, bound, or physically forced. But violence cannot directly force the will itself. A tyrant may force a person's hand, but he cannot make the person inwardly consent in the same way he can move the body.

Fear creates mixed cases. A person may throw cargo overboard in a storm or obey under threat. Such actions are not what the person would choose in themselves, but under the circumstances they may still be chosen as a means to avoid a greater evil. Fear can reduce blame, but it does not automatically erase responsibility.

Desire, or concupiscence, also does not automatically make action involuntary. In fact, strong desire often makes an act more voluntary, because the person moves eagerly toward the desired object. Desire may cloud reason, but it usually reveals what the will is pursuing.

Ignorance is different. Ignorance can make an act involuntary when the person does not know an important fact and would not have acted that way if he had known. But ignorance can also be blameworthy if the person should have known better. Aquinas keeps the same hard distinction students have seen in Aristotle: not all ignorance excuses.

\subsection*{6. Habit: Freedom Becomes Stable or Unstable}
Aquinas gives special attention to habit. A habit is a stable disposition by which a person is ready to act well or badly. Habits are not passing moods. They are formed by repeated acts and make future acts easier.

Aquinas follows Aristotle here: by repeated actions we become disposed toward similar actions. A person who repeatedly tells the truth becomes more ready to be truthful. A person who repeatedly gives in to anger becomes more ready to answer with anger. A person who practices courage becomes more able to stand firm when fear arrives.

This means freedom is not only a moment of choosing. Freedom is formed. Earlier choices shape what later choices feel like. A person may become freer for the good by building virtue, or less free by building vice.

Aquinas's view is demanding but hopeful. Bad habits can bind a person, but habits are not fate. Because human beings have reason and will, they can be taught, corrected, disciplined, and re-formed toward better action.

\subsection*{7. Law Is Not Mere Force}
Aquinas then turns to law. Many students may think law is simply a rule backed by punishment. Aquinas gives a richer definition: law is an ordinance of reason for the common good, made by one who has care of the community, and made known.

This definition matters. Law is not mere will. If a ruler simply commands whatever pleases him, that is closer to lawlessness than law. Real law belongs to reason. It directs human action toward a good that is more than private advantage.

Law is also for the common good. It should not exist merely for the ruler, the strong, or one private group. It should guide the community toward justice, peace, and human flourishing.

Finally, law must be made known. A hidden rule cannot guide action well. To bind rational creatures, law must be promulgated; it must be given in a way that can actually direct reason and choice.

\subsection*{8. Natural Law and Human Freedom}
Aquinas says that human beings participate in eternal law through natural law. Eternal law is God's wise ordering of all things. Natural law is the way rational creatures share in that order by reason.

The first principle of natural law is simple: good is to be done and pursued, and evil is to be avoided. Other basic goods follow from human nature: preserving life, forming families, educating children, seeking truth, living in society, avoiding ignorance, and treating others in ways that allow common life.

Natural law does not mean that every moral question is easy. Human beings can make mistakes, and passions, customs, bad habits, and corrupt societies can distort judgment. But Aquinas thinks reason is not blind. We have some real light by which we can recognize basic goods.

This again changes the meaning of freedom. Freedom is not the power to invent good and evil from nothing. Freedom is the power of a rational creature to recognize the good, choose means to it, and be formed by law and habit toward a life worthy of a human being.

\subsection*{9. Is Law an Enemy of Freedom?}
For Aquinas, good law does not destroy freedom. Good law teaches, directs, restrains, and forms. It is like a rule of reason that helps human beings become capable of the good.

A child may first experience law as outside control: do this, do not do that. But if the law is reasonable and good, it trains the person to see and love what is good. Over time, a virtuous person does not merely obey from fear. He acts from an inwardly formed love of the good.

This is why Aquinas can connect law and freedom. A person ruled by disordered desire may feel free because no one is stopping him, but he may actually be enslaved to appetite, pride, fear, or habit. A person formed by reason, virtue, and good law may be more truly free because he can choose the good steadily.

Aquinas therefore stands between Augustine and the later compatibilists. With Augustine, he thinks the will must be ordered toward higher goods. With Aristotle, he thinks habit and reason shape responsibility. And against a shallow idea of freedom, he argues that freedom is completed, not destroyed, by reason's order toward the good.

\section*{Reading Focus}
\begin{enumerate}[leftmargin=*]
\item What is Aquinas's distinction between a \textit{human act} and an \textit{act of a human}?
\item Why does Aquinas think human actions are done for an end?
\item How can outside influences exist without automatically destroying voluntary action?
\item Why does Aquinas think habit matters for freedom and responsibility?
\item According to Aquinas, why is good law not merely force or control?
\end{enumerate}

\section*{Reason--Will--End Map and Summary}
Create a map showing how Aquinas connects \textbf{reason}, \textbf{will}, \textbf{end}, \textbf{habit}, and \textbf{law}. Include at least one drawing or symbol for each part. At the bottom, write a 3--5 sentence summary explaining whether Aquinas thinks law is an enemy of freedom or a teacher of freedom.

\section*{Connection to the Unit Question}
Write one claim answering this question: \textbf{Is freedom the ability to choose anything, or the ability to choose the good rationally?} Use Aquinas to support your answer, and compare him to either Aristotle or Augustine.
\end{multicols}
\end{document}
